% ============================================================
% conclusion_draft.tex
% ============================================================

\newpage
\section{Conclusion}

This thesis tested whether CNN-based chart pattern models have predictive power for short-term equity returns in the Euro Stoxx 50, and whether more advanced architectures, better image encodings, and hybrid designs improve on the baseline model from \textcite{Jiang2023}.

\subsection{Summary of Findings}

\paragraph{H1: Predictive Power in European Equities.}
All three models achieve out-of-sample AUC above 0.5 with confidence intervals entirely above the 0.5 baseline. The Baseline CNN reaches 0.5166, EfficientNet-B2 reaches 0.5160, and ConvNeXt V2 reaches 0.5179. The signal is small but consistent across models and confirmed by positive IC values. A corrected portfolio backtest using actual 20-day forward returns produces positive gross annualised returns of approximately 5\% for the Baseline CNN and ConvNeXt V2, though the 65 non-overlapping test periods are too few for the portfolio returns to reach conventional statistical significance. $H_{10}$ is rejected: CNN-based chart models do carry a predictive signal in European equities.

\paragraph{H2: Architecture and Image Encoding.}
The advanced hybrid architectures do not significantly outperform the Baseline CNN in terms of AUC. The differences are small (below 0.005 AUC points), fall within single-run training variance, and are not significant at conventional levels. The image encoding comparison experiment reveals a statistically significant result: the v4 configuration consistently outperforms the v3 configuration for both EfficientNet and ConvNeXt ($p < 0.01$ in both cases), while within each image format the two architectures perform at similar levels. Across the configurations tested, the input design choice matters more than the backbone choice.

One caveat applies to the encoding comparison: the v3 and v4 configurations also differ in the number of technical indicators supplied to the sequence encoder (21 for v3, 25 for v4), so the performance advantage cannot be attributed solely to the image format without a controlled ablation. The conclusion is that the v4 configuration as a whole is more effective, with the multi-scale image structure as the most likely primary driver.

This points to image design as an underexplored dimension in the chart pattern CNN literature. Most prior work focuses on architectural depth, regularisation, or pretraining, while keeping the image format fixed as a single-scale greyscale chart. The multi-scale RGB approach provides additional information that the models can exploit, particularly in trending markets.

\paragraph{H3: Geographical Generalisability.}
European AUC values of 0.516--0.518 fall somewhat below the 0.54--0.56 reported for U.S. equities in \textcite{Jiang2023}, but a formal statistical comparison is not possible: the evaluation periods do not overlap, the universes differ by an order of magnitude in size, and the post-pandemic period may be structurally harder to predict than the pre-2020 period. H3 is therefore treated as a descriptive range check rather than a hypothesis test. Taken at face value, the European results are in the same general order of magnitude as the U.S. benchmark, suggesting that chart-based CNN signals are not confined to the U.S. market, while also indicating that European predictability may be somewhat lower.

\paragraph{H4: Cryptocurrency Transfer.}
ConvNeXt V2 achieves an AUC of 0.5433 on BTC-USD and ETH-USD without any retraining, compared to its equity AUC of 0.5179. The model was trained entirely on Euro Stoxx 50 equities and shows better discrimination on cryptocurrency data, a result that supports the hypothesis that the learned visual features capture market-agnostic chart patterns rather than equity-specific regularities. $H_{40}$ is rejected: the equity-trained model does retain meaningful predictive power when applied to cryptocurrency markets.

\subsection{Implications}

The practical implication of the image encoding finding is that researchers and practitioners building chart pattern models should treat the image design as a hyperparameter to be tuned. Multi-scale representations that encode price dynamics at multiple time horizons appear to carry more information than single-scale images, regardless of how sophisticated the underlying model is. This is a relatively cheap improvement: the v4 format requires no additional data and only modestly more preprocessing than a standard OHLC chart.

The crypto transfer result is relevant for applications where training data is scarce. Cryptocurrency markets have shorter histories than equity markets, and directly training large models on a few years of daily data is difficult. The fact that a model trained on 19 years of European equity data generalises well to crypto suggests that transfer learning across asset classes may be viable for chart-based models, at least at the level of daily price patterns.

The portfolio backtest, conducted on actual 20-day forward returns with non-overlapping rebalancing, shows that Baseline CNN and ConvNeXt V2 produce positive gross Sharpe Ratios of around 0.4--0.5 from a quintile long-short strategy. These are not statistically significant given the 65 independent test periods, but they are directionally consistent with the AUC finding. Whether the signal is large enough to be economically exploitable is uncertain: the equal-weighted benchmark achieved a Sharpe of 1.04 over the same period, and the long-short spread does not approach that level. Translating the modest predictive signal identified by AUC into a competitive trading strategy would likely require a substantially larger asset universe to reduce idiosyncratic noise per portfolio leg.

\subsection{Suggestions for Future Research}

Several natural extensions follow from this work.

First, the image encoding comparison covered only two formats. A systematic study of image design choices, including different panel layouts, colour mappings, resolution, and temporal scale combinations, would help identify which visual features are most useful for the models. The finding that v4 outperforms v3 is a single data point in what could be a much larger design space.

Second, this thesis evaluated a fixed single train-validation-test split. Rolling window re-training, where the model is retrained periodically as new data arrives, could reduce the mismatch between training conditions and the out-of-sample period and may improve temporal stability.

Third, the crypto transfer result merits a more careful follow-up with a broader set of assets and an explicitly designed cross-asset training strategy. Fine-tuning an equity-pretrained model on a small amount of crypto data, rather than zero-shot transfer, would likely improve performance further.

Fourth, combining chart images with other data types, such as earnings announcements, short interest, or options market signals, could be a productive direction. The hybrid architecture already supports multiple inputs through the fusion module and could be extended without major changes to the overall design.

Finally, the robustness analysis in Section~\ref{sec:res_stoxx600} already extends the classification comparison to the full STOXX Europe 600, and the small but real signal survives on that much larger and more heterogeneous cross-section. Two questions remain open. The first is economic: the portfolio backtest was run only on the 50-stock universe, and repeating it on the 600-stock cross-section, where each long-short leg would hold far more names and idiosyncratic noise would diversify away more effectively, would test whether the broader universe is enough to turn the directional signal into an economically meaningful strategy. The second is reach: even the STOXX Europe 600 is restricted to established large-, mid-, and small-cap constituents, so extending the analysis to micro-cap and less liquid European names would clarify how far down the liquidity spectrum the signal persists.
